\section{Chuỗi suy luận là hành vi quan sát được}
\label{sec:ch02-chuoi-suy-luan}

\index{chuỗi suy luận}%
Các mục trước đã cho thấy đầu ra phụ thuộc cả vào trọng số lẫn điều kiện sử
dụng. \tn{Chuỗi suy luận}{chain-of-thought} đặt sự phụ thuộc đó ngay trong
văn bản đầu ra: \tn{lời nhắc}{prompt} chứa các ví dụ có bước trung gian để gợi mô hình sinh
một chuỗi bước trước khi nêu đáp án. Wei và cộng sự báo cáo rằng cách nhắc
này có thể cải thiện các bài toán suy luận nhiều bước trên những mô hình đủ
lớn~\autocite{p06cot}.

Tuy nhiên, một chuỗi bước được sinh ra vẫn trước hết là hành vi cần quan sát.
Nó cho thấy mô hình đã in những token nào dưới lời nhắc đã chọn;
nó không tự cho thấy token nào trong chuỗi là nguyên nhân của đáp án, hay
chuỗi ấy có khớp với phép tính tạo ra đáp án không. Ngay bài báo gốc cũng ghi rằng không có gì bảo đảm đường suy luận là đúng,
nên một chuỗi được sinh ra có thể dẫn tới đáp án đúng lẫn đáp án
sai~\autocite{p06cot}.

Anand và cộng sự, trong khảo sát về suy luận của các mô hình này, để ngỏ
chính câu hỏi liệu suy luận có phải là một năng lực đơn nhất hay là hợp thành
của nhiều quá trình tách rời được. Thay vì trả lời, khảo sát mô tả suy luận
như một hành vi thay đổi theo điều kiện đo, và liệt kê các điều kiện ấy: quy
mô mô hình, cấu trúc nhiệm vụ, lời nhắc, công cụ hoặc bộ nhớ ngoài, giám sát
và giao thức đánh giá~\autocite{p08periodic}. Cuốn sách mượn đúng cách nhìn
ấy, rằng suy luận là hành vi đo được dưới một điều kiện chứ không phải một
năng lực có sẵn. Khi một hệ thống trả lời tốt hơn sau khi được yêu cầu trình
bày bước trung gian, ta có bằng chứng về một thay đổi hành vi dưới lời nhắc
ấy. Ta chưa có bằng chứng rằng văn bản trung gian là lời giải thích trung thực
về cơ chế bên trong.

